\Section{Символика и примеры её использования}

Заглавные латинские буквы \(A\), \(B\), \(C\) и т.д.~--- символы множеств.

Прописные латинские буквы \(a\), \(b\), \(c\) и т.д.~--- символы элементов
множеств.

Иногда для обозначения числовых множеств будут использоваться греческие буквы
\(\alpha\), \(\beta\), \(\gamma\) и т.д.

\(a \in A\) ~--- <<\(a\) является элементом множества \(A\)>> или
<<\(a\) принадлежит множеству \(A\)>>.

\(a \notin A\) ~--- <<\(a\) не является элементом множества \(A\)>> или
<<\(a\) не принадлежит множеству \(A\)>>.

\(A \subset B\) ~--- <<\(A\) является подмножеством множества \(B\)>>.

\(A \not\subset B\) ~--- <<\(A\) не является подмножеством множества \(B\)>>.

\(A \cup B\) ~--- <<объединение \(A\) и \(B\)>>.

\(A \cap B\) ~--- <<пересечение \(A\) и \(B\)>>.

\(A \setminus B\) ~--- <<разность \(A\) и \(B\)>>.

\(\forall\) ~--- квантор всеобщности (<<для любого>>, <<для всех>>).

\(\exists\) ~--- квантор существования (<<существует>>).

\(\lor\) ~--- дизъюнкция (<<или>>).

\(\&\) ~--- конъюнкция (<<и>>).

\(\Rightarrow\) ~--- импликация (<<следует>>, <<то>>, <<тогда>>).

\(\Leftrightarrow\) ~--- эквивалентность (<<тогда и только тогда>>,
<<необходимо и достаточно>>).

\(\neg\) ~--- отрицание (<<не>>).

\(\mapsto\) ~--- <<выполняется>>, <<имеет место>>.

\(\eqdef\) ~--- равенство по определению (<<по определению>>).

\(=\left\{x:\right\}\) ~--- множество, которое описывается через элементы

\(:\) ~--- <<такой, что>>.

\begin{example}
    Мн-во \(A\) явл. подмн-вом мн-ва \(B\), если любой элемент мн-ва \(A\) явл.
    элементом мн-ва \(B\):

    \[
        \left[A \subset B\right] \eqdef
        \left[\forall a \in A \mapsto a \in B\right]
    \]
\end{example}

\begin{example}
    Мн-ва \(A\) и \(B\) равны, если \(A\) явл. подмн-вом \(B\) \underline{и}
    \(B\) явл. подмн-вом \(A\):

    \begin{gather}
        \left[A = B\right] \eqdef \left[A \subset B\right] \&
        \left[B \subset A\right] \\
        \left[A = B\right]\eqdef \left[\forall a \in A \mapsto a \in B\right]
        \& \left[\forall b \in B \mapsto b \in B\right]
    \end{gather}
\end{example}

\begin{example}
    Объединением мн-в \(A\) и \(B\) наз. мн-во \(A \cup B\), каждый элемент
    которого принадлежит \(A\) \textbf{или} \(B\):

    \begin{gather}
        A \cup B = \left\{x: [x \in A] \lor [x \in B]\right\} \\
        \left[x \in A \cup  B\right] = \left[x \in A\right] \lor \left[x \in
        B\right]
    \end{gather}

\end{example}

\begin{example}
    Пересечние множеств \(A\) и \(B\) называется множество \(A \cap B\), каждый 
    элемент которого принадлежит \(A\) \textbf{и} \(B\):

    \begin{gather}
        A \cap B = \left\{x: [x \in A] \& [x \in B]\right\} \\
        \left[x \in A \cap  B\right] = \left[x \in A\right] \& \left[x \in
        B\right]
    \end{gather}
\end{example}

\begin{example}
    Разностью множеств \(A\) и \(B\) называется множество \(A \setminus B\),
    каждый элемент которого принадлежит \(A\) и не принадлежит \(B\):

    \begin{gather}
        A \setminus B = \left\{x: [x \in A] \lor [x \notin B]\right\} \\
        \left[x \in A \cup  B\right] = \left[x \in A\right] \lor \left[x \notin
        B\right]
    \end{gather} 
\end{example}

Теперь начнем формулировать отрицание высказываний.

\begin{example}
    При отрицании, квантор всеобщности \(\forall\) заменяется на квантор
    существования \(\exists\). Так же после квантора всеобщности \(\forall\)
    принято ставить <<\(\mapsto\)>>, а после квантора существования \(\exists\)
    ~--- <<:>> (это правило не является строгим, символы стоит использовать в
    зависимости от их произношения).

    Обозначим \(\mathfrak{A} = \left[A \subset B\right] \eqdef
    \left[\forall a \in A \mapsto a \in B\right] \), тогда:

    \[
    \neg \mathfrak{A} \eqdef \left[A \not\subset B\right] \eqdef
    \left[\exists a \in A: a \notin B\right]
    \]

\end{example}

\begin{example}
    При отрицании, все символы, которые находятся при кванторах
    (\(\forall\), \(\exists\)) не меняются, а все остальные ~---
    заменяются на противоположные (\(\in\) ~--- на \(\notin\),
    \(>\) ~--- на \(\leq\), \(\lor\) ~--- на \(\&\) и т.д.).

    Обозначим \( \mathfrak{B} =
    \left[\forall x \in X \mapsto \lvert x \rvert < \alpha \right] \), тогда:

    \[
    \neg \mathfrak{B} = \left[\exists x \in X: \lvert x \rvert
    \ge \alpha \right]
    \]
\end{example}

\begin{example}
    Теперь рассмотрим пример с несколькими кванторами

    Обозначим \(\mathfrak{C} = \left[\exists \alpha > 0: \forall x \in X
    \mapsto \lvert x \rvert < \alpha\right]\), тогда:
    \[
    \neg \mathfrak{C} = \left[ \forall \alpha > 0\;\exists x_\alpha \in X:
    \lvert x_\alpha \rvert >= \alpha\right]
    \]

    Стоит обратить внимание на то, что при отрицании мы добавили к \(x\) индекс
    \(\alpha\). Это связано с тем, что для каждого \(\alpha\) будет существовать
    свой \(x_{\alpha}\), так что использование \(x\) в данном случае является
    некорректным.
\end{example}

Стоит так же отметить тот факт, что запись через символы является довольно
свободной, и у нее нет строгих правил записи.

\Section{Теория действительного числа}
\Subsection{Рациональные числа и их свойства}

\begin{definition}
    Рациональным числом называется дробь \(\frac{p}{q}\), где
    \(q \in \mathbb{N}\), \(p \in \mathbb{Z}\).
\end{definition}

Упорядоченность \(\mathbb{Q}\):

Для любых рациональных чисел \(a\) и \(b\) имеет место только одно из
соотношений: \(a = b\), \(a > b\), \(b > a\).

Запишем основные свойства рациональных чисел:

\begin{enumerate}
    \item Если \(a = b\) и \(b = c\), то \(a = c\),

    Если \(a > b\) и \(b > c\), то \(a > c\).

    (транзитивность <<\(=\)>> и <<\(>\)>>)
\end{enumerate}

\textbf{Сложение рациональных чисел:}

\(a = \frac{p}{q}\) и \(b = \frac{r}{s}\).

Сумма \(a + b = \frac{ps + rq}{qs}\).

Операция нахождения суммы называется сложением.

\begin{enumerate}[resume]
    \item \(a + b = b + a\) (коммутативность сложения);
    \item \((a + b) + c = a + (b + c)\) (ассоциотивность сложения);
    \item Существует число \(0\) такое, что для любого \(a\) выполняется
    \(a + 0 = a\) (особая роль нуля);
    \item Для любого числа \(a\) найдется противоположное число \(-a\) такое,
    что \(a + (-a) = 0\) (существование противоположных чисел);
\end{enumerate}

\textbf{Умножение рациональных чисел}

\(a = \frac{p}{q}\) и \(b = \frac{r}{s}\).

Произведение \(a \cdot b = \frac{pr}{qs}\).

Операция нахождения произведения называется умножением.

\begin{enumerate}[resume]
    \item \(a \cdot b = b \cdot a\) (коммутативность умножения);
    \item \((a \cdot b) \cdot c = a \cdot (b \cdot c)\) (ассоциативность
    умножения);
    \item Существует число 1 такое, что для любого \(a\) выполняется
    \(a \cdot 1 = a\) (особая роль 1);
    \item Для любого \(a \not = 0\) существует обратноечисло \(\frac{1}{a}\)
    такое, что \(a \cdot \frac{1}{a} = 1\) (существование обратных чисел);
    \item \((a + b) \cdot c = a \cdot c + b \cdot c\)
    (дистрибутивность сложения относительно умножения);
    \item Если \(a > b\), то для любого \(c\) выполняется \(a + c > b + c\);
    \item Если \(a > b\), то для любого \(c > 0\) выполняется
    \(a \cdot c > b \cdot c\);
    \item Для любого \(a > 0\) найдется \(n \in \mathbb{N}\) такое, что
    \(n > a\) (аксиома Архимеда).
\end{enumerate}

\textbf{Плотность рациональных чисел}

Если \(a > b\), то найдется \(c \in \mathbb{Q}\) такое, что \(a > c > b\)

Для доказательства этого утверждения достаточно взять \(c = \frac{a + b}{2}\).